\chapter{Introdução}
Apresente aqui o contexto, o problema de pesquisa, a justificativa e a organização do trabalho. As diretrizes gerais de estrutura e apresentação de trabalhos acadêmicos estão definidas na norma vigente \cite{abnt14724}; um exemplo de citação no texto pode ser feito com \citeonline{gil2022}.

A Figura~\ref{fig:exemplo} ilustra a inserção de uma figura baseada no modelo institucional~\cite{depro2021}, o Quadro~\ref{qua:exemplo} exemplifica um quadro editável (fechado nas quatro laterais, voltado a dados qualitativos) e a Tabela~\ref{tab:exemplo} demonstra uma tabela numérica no padrão tabular aberto.

\begin{figure}[htbp]
  \centering
  \caption{Título da figura}\label{fig:exemplo}
  \includegraphics[height=1.8cm,alt={Logotipo do Departamento de Engenharia de Produção da UFS utilizado como figura de exemplo}]{figuras/logo-depro.pdf}
  \fonte{\citeonline{depro2021}}
\end{figure}

\begin{quadro}[htbp]
  \centering
  \caption{Título do quadro}\label{qua:exemplo}
  \tagpdfsetup{table/header-rows={1}}
  \begin{tabularx}{\textwidth}{|>{\bfseries\raggedright\arraybackslash}p{3.8cm}|>{\raggedright\arraybackslash}X|>{\raggedright\arraybackslash}X|}
    \hline
    \rowcolor{gray!12}
    \textbf{Coluna A} & \textbf{Coluna B} & \textbf{Coluna C} \\ \hline
    Exemplo & Texto editável & Resultado \\ \hline
    Outro exemplo & Texto editável & Resultado \\ \hline
  \end{tabularx}
  \fonte{Elaboração própria}
\end{quadro}

\begin{table}[htbp]
  \centering
  \caption{Título da tabela}\label{tab:exemplo}
  \tagpdfsetup{table/header-rows={1}}
  \begin{tabular}{lrrr}
    \toprule
    \textbf{Item} & \textbf{Quantidade (unid.)} & \textbf{Área ($\text{m}^2$)} & \textbf{Custo (R\$)} \\
    \midrule
    A & 6.529 & 9.482 & 647 \\
    B & 4.894 & 6.524 & 644 \\
    C & 556 & 5.675 & 238 \\
    D & 2.546 & 649 & 787 \\
    \midrule
    \textbf{Total} & \textbf{14.525} & \textbf{22.330} & \textbf{2.316} \\
    \bottomrule
  \end{tabular}
  \fonte{Elaboração própria}
\end{table}

Um exemplo de equação numerada e referenciada como Equação~\eqref{eq:area} é dado por:
\begin{equation}\label{eq:area}
  A = \pi r^2
\end{equation}
\needspace{4\baselineskip}
\noindent em que:
\par\nopagebreak\vspace{0.3em}\nopagebreak
\begin{description}
  \item[$A$] área de um círculo, medida em $[\text{m}^2]$;
  \item[$r$] raio da circunferência.
\end{description}
